% TOP_PROBLEM: {"id": 3000039, "title": "Goddyn's conjecture on thin spanning trees", "category_id": 2, "category": {"id": 2, "name": "combinatorics", "display_name": "Combinatorics", "description": "Counting problems, graph theory, discrete structures.", "slug": "combinatorics", "order_index": 2, "created_at": "2026-07-31T15:26:25.671Z"}, "status": "partially_solved", "problem_number": "AMR-029-0039", "set_id": 13}
% FIRST_POSTED: 2026-10-06
% ATTEMPT_STATUS: solved
% OPENAI_CLAIM: {"schema_version":1,"source_repo":"https://github.com/openai/math","source_commit":"adc7f1241b42e322a6451854ab7e4b4c146bf78a","release_date":"2026-10-06","match":"stronger","resolution":"proved","independently_reviewed":false,"second_pass":"agreed","families":[{"family":"174","title":"Deterministic construction of strong thin spanning trees","subject":"Combinatorics","manuscripts":[{"title":"A polynomial-time construction of strong thin trees","pdf_path":"preprints/A-polynomial-time-construction-of-strong-thin-trees-September-23-2026/paper.pdf","readme_path":"preprints/A-polynomial-time-construction-of-strong-thin-trees-September-23-2026/README.md","date":"2026-09-23","theorem_ref":"Theorem 1.1"},{"title":"The strong thin tree conjecture","pdf_path":"preprints/The-strong-thin-tree-conjecture-September-23-2026/paper.pdf","readme_path":"preprints/The-strong-thin-tree-conjecture-September-23-2026/README.md","date":"2026-09-23","theorem_ref":"Theorem 1.1"}],"lean":{"doc_path":"lean/docs/174.md","comparators":["lean/ComparatorChallenges/AlgorithmicThinTrees.json","lean/ComparatorChallenges/AlgorithmicThinTrees.lean","lean/ComparatorChallenges/StrongThinTree.json","lean/ComparatorChallenges/StrongThinTree.lean"],"declarations":["OAI.StrongThinTree.strongThinTree"],"covers_main_theorem":true},"reasoning_trace":null}]}
\documentclass[11pt]{article}
\usepackage[margin=25mm]{geometry}
\usepackage{fontspec}
\setmainfont{Latin Modern Roman}
\usepackage{amsmath,amssymb,mathtools}
\usepackage{unicode-math}
\setmathfont{Latin Modern Math}
\usepackage{enumitem}
\usepackage{xurl,hyperref,bookmark}
\hypersetup{colorlinks=true,urlcolor=blue}
\setcounter{secnumdepth}{0}
\setlength{\parindent}{0pt}
\setlength{\parskip}{5pt}
\setlength{\emergencystretch}{3em}
\begin{document}
\section*{Goddyn's conjecture on thin spanning trees}

\subsection{OpenAI's claim}
\textit{Algorithmic strong thin trees}: The algorithmic manuscript claims a deterministic procedure polynomial in the binary input length that returns a $C/k$-thin spanning tree in a finite $k$-edge-connected loopless multigraph; binary-encoded parallel multiplicities and the one-vertex case are included.

\textit{Existence of strong thin trees}: The existence paper claims one universal $C$ such that a $k$-edge-connected loopless finite multigraph has a spanning tree using at most a $C/k$ fraction of each cut; this assertion is simultaneous over all cuts of the graph.

\subsection{Outline of the argument}
\textit{Algorithmic strong thin trees}: An effective metric-cone hierarchy constructs a rational shortcut form and many low-resistance candidate edges. A second frame preserves the ranks needed for spanning-tree packing, and rational padding supplies signing certificates without irrational whitening. Iterated simultaneous sparsification controls every cut and packing inequality, with a separate reduction handling binary multiplicities.

\textit{Existence of strong thin trees}: A shortcut quadratic form supplies a cut-sensitive representation of suitable edges, while an isotropic representation records the partition inequalities for tree packing. Rank-one signing selects the same edge subset in both representations. Iterated halving has a bounded product of relative losses, so the remaining tree retains uniform cut control.

\subsection{Scope relative to this problem}
\textit{Algorithmic strong thin trees}: The uniform C/k cut ratio allows f(epsilon)=ceil(C/epsilon), proving Goddyn’s existential connectivity-threshold assertion with an explicit linear dependence; no simplicity hypothesis is needed.

\textit{Existence of strong thin trees}: The uniform C/k cut ratio allows f(epsilon)=ceil(C/epsilon), proving Goddyn’s existential connectivity-threshold assertion with an explicit linear dependence; no simplicity hypothesis is needed.

The claimed result implies the stated question and also supplies the additional strength described above.

\subsection{Formal verification}
Family 174 describes Lean coverage for thin-tree existence and the companion deterministic construction with binary multiplicities.

\textit{Algorithmic strong thin trees}: The manuscript links \texttt{AlgorithmicThinTrees.lean}, \texttt{StrongThinTree.lean} (declaration \texttt{OAI.StrongThinTree.strongThinTree}).

The documented selection includes this manuscript’s main assertion; that assertion retains the restrictions stated above.

\textit{Existence of strong thin trees}: The manuscript links \texttt{AlgorithmicThinTrees.lean}, \texttt{StrongThinTree.lean} (declaration \texttt{OAI.StrongThinTree.strongThinTree}).

The documented selection includes this manuscript’s main assertion; that assertion retains the restrictions stated above.

This site has not run the Lean development. These links describe OpenAI’s selected formal statements and should be read with their stated scope.

\subsection{Status caveat}
OpenAI reports that an internal model produced the manuscripts. Their arguments have not been independently verified as proofs by this site. Any portion outside the documented formal scope may contain mistakes, and the catalogue label reports the announced mathematical claim.

\subsection{Sources}
\begin{itemize}
\item OpenAI, \emph{A polynomial-time construction of strong thin trees}, 2026-09-23; citation key \texttt{OAI:A-polynomial-time-construction-of-strong-thin-trees-September-23-2026}. \href{https://github.com/openai/math/blob/adc7f1241b42e322a6451854ab7e4b4c146bf78a/preprints/A-polynomial-time-construction-of-strong-thin-trees-September-23-2026/paper.pdf}{PDF (GitHub)}; \href{https://github.com/openai/math/raw/adc7f1241b42e322a6451854ab7e4b4c146bf78a/preprints/A-polynomial-time-construction-of-strong-thin-trees-September-23-2026/paper.pdf}{PDF (direct)}; \href{https://github.com/openai/math/blob/adc7f1241b42e322a6451854ab7e4b4c146bf78a/preprints/A-polynomial-time-construction-of-strong-thin-trees-September-23-2026/README.md}{Manuscript page and citation}; \href{https://github.com/openai/math/blob/main/preprints/A-polynomial-time-construction-of-strong-thin-trees-September-23-2026/paper.pdf}{Latest version}.
\item OpenAI, \emph{The strong thin tree conjecture}, 2026-09-23; citation key \texttt{OAI:The-strong-thin-tree-conjecture-September-23-2026}. \href{https://github.com/openai/math/blob/adc7f1241b42e322a6451854ab7e4b4c146bf78a/preprints/The-strong-thin-tree-conjecture-September-23-2026/paper.pdf}{PDF (GitHub)}; \href{https://github.com/openai/math/raw/adc7f1241b42e322a6451854ab7e4b4c146bf78a/preprints/The-strong-thin-tree-conjecture-September-23-2026/paper.pdf}{PDF (direct)}; \href{https://github.com/openai/math/blob/adc7f1241b42e322a6451854ab7e4b4c146bf78a/preprints/The-strong-thin-tree-conjecture-September-23-2026/README.md}{Manuscript page and citation}; \href{https://github.com/openai/math/blob/main/preprints/The-strong-thin-tree-conjecture-September-23-2026/paper.pdf}{Latest version}.
\item Family 174: \href{https://github.com/openai/math/blob/adc7f1241b42e322a6451854ab7e4b4c146bf78a/lean/docs/174.md}{Lean scope}.
\item \href{https://github.com/openai/math/blob/adc7f1241b42e322a6451854ab7e4b4c146bf78a/lean/ComparatorChallenges/AlgorithmicThinTrees.json}{Comparator AlgorithmicThinTrees.json}.
\item \href{https://github.com/openai/math/blob/adc7f1241b42e322a6451854ab7e4b4c146bf78a/lean/ComparatorChallenges/AlgorithmicThinTrees.lean}{Comparator AlgorithmicThinTrees.lean}.
\item \href{https://github.com/openai/math/blob/adc7f1241b42e322a6451854ab7e4b4c146bf78a/lean/ComparatorChallenges/StrongThinTree.json}{Comparator StrongThinTree.json}.
\item \href{https://github.com/openai/math/blob/adc7f1241b42e322a6451854ab7e4b4c146bf78a/lean/ComparatorChallenges/StrongThinTree.lean}{Comparator StrongThinTree.lean}.
\item \href{https://github.com/openai/math/blob/adc7f1241b42e322a6451854ab7e4b4c146bf78a/CONTENTS.md}{OpenAI catalogue entry}.
\end{itemize}
\end{document}
